% Main figures: files and legends (Gate C). Each legend's title is the figure's title (its build docstring) and its claim;
% the body says what each panel plots, its n, its summary marks, its test and correction, and the scale bars. The facts
% come from the figure source data (manuscript/writing/legend_facts.py); the main legends count toward the 7,000 words.
\setfigfile{1}{../figures/v4/main/fig1.pdf}
\setfigfile{2}{../figures/v4/main/fig2.pdf}
\setfigfile{3}{../figures/v4/main/fig3.pdf}
\setfigfile{4}{../figures/v4/main/fig4.pdf}
\setfigfile{5}{../figures/v4/main/fig5.pdf}
\setfigfile{6}{../figures/v4/main/fig6.pdf}
\setfigfile{7}{../figures/v4/main/fig7.pdf}

\deflegend{1}{A whole-CNS connectome learns to plan and to remember}{%
(A) The unfolded CNS of the model: every located cell body, colored by superclass group. Scale bar, 100~μm.
(B) The interface: the visual input cells (white), the arbors of the \val{L04} readout cells and the giant fiber DNp01; insets,
an input cell (LC12) and a readout cell (i2 MN). Scale bars, 100~μm; insets, 20~μm (LC12) and 50~μm (i2 MN).
(C and D) One memory trial (scene 1, readout trained, cue 0): the cells that spike in the cue, the delay and the input (C), and
the cells spiking per ms by group (D).
(E) The model at a glance: input and readout pathways, fixed and trained parameters, and the readout cells by their shortest
route from the input.
(F) The tasks: three of the \val{L09} planning maps with their optimal moves, and the delayed-cue protocol with its \val{A01}
trained of \val{A04} timings.
(G) Learning cost against certified difficulty, \val{L09} maps in \val{L10} replicate sets; lines, least-squares fits on log
axes; inset, Spearman ρ per set with permutation tests.
(H) Maps solved by each pairing of brain and readout checkpoints; McNemar tests, BH-adjusted.
(I) Net memory use at trained (filled) and untrained (open) timings in \val{L15} replicates of four learners (\val{L19},
\val{L19}, \val{L20} and \val{L20}), joined by replicate; bars, medians; exact sign-flip tests of the lock (trained minus
untrained use), Holm-adjusted.
See also Figure~S1 and Video~S1.}

\deflegend{2}{Working memory is held in an activity-silent intrinsic trace}{%
(A) One trial (scene 1, readout trained, \val{G06}-ms delay), ms by ms: cells with a voltage difference between the cues (log
scale) and cells holding the adaptation trace, by group.
(B) Cells that spike only after cue 0 or only after cue 1, with the cells per ms above; shaded, none.
(C) The trace (adaptation, cue 1 minus cue 0) at five moments of the trial; white, cells spiking differently. Scale bar,
100~μm.
(D) Decay of the cue difference by state variable, relative to early in the delay; lines, exponential fits (time constants
over four learners).
(E) The share of the voltage difference from adaptation at the end of the delay and from synaptic input at the first decision,
four learners, after the \val{G06}-ms (filled) and \val{H12}-ms (open) delay.
(F) Transplants of the full state or of adaptation alone at the end of the delay: cosine with the native decision state.
(G) Trained histories solved with the own or the other cue's state, all trained timings, four learners; exact tests,
Holm-adjusted.
(H) Top, the synaptic share of the decision-time drive at all \val{H07} timings, four learners; dashed, the minimum. Bottom,
the cue drive into readout types from types that held the trace, by region; dashed, their share of types.
(I) Cumulative differential spikes through the \val{H12}-ms delay, four learners.
(J) The trace at the decision by delay and prefix input; the line without input is the fitted decay.
(K) The cue drive at the first decision between cell types of the brain (above) and ventral nerve cord (below): width, drive;
color and arrowhead, the sender's transmitter; gold, types that held the trace; outlined, readout types.
See also Figure~S2 and Video~S2.}

\deflegend{3}{The trace is distributed, and no circuit is special}{%
(A) The scene-1 cue ensemble (the \val{D15} of cells with the largest cue difference at the decision) beside three of its
\val{D14} matched sets, which share its cell types, distance from the input, degree and readout role; right, four cue-tuned
cell types. Scale bars, 100~μm and 50~μm (cell types).
(B) The ensemble's connections as a matrix (sources down, targets across) with its densest block enlarged, and its internal
connections by source group against the mean of its \val{D14} matched sets, two scenes.
(C) Silencing the ensemble at the decision: the change in decision margin relative to the matched sets' mean (gray, each
matched set), ensembles of two scenes tested in both; one-sided tests, Holm-adjusted.
(D) Removing the ensemble's internal synapses, against matched synapses, as in (C).
(E) Whole populations removed: trained histories solved of \val{D07}; dashed, healthy; equivalence tests (TOST),
Holm-adjusted.
(F) Maps of the cue difference at the decision, and each circuit's share of the cue against its share of cells, three scenes.
(G) Where three connection rules act, and trained histories solved under each variant rule (three runs each), healthy
(filled) and under the pan-neuronal Kv4 lesion (open).
See also Figure~S3.}

\deflegend{4}{Trace-carrying cells have a molecular signature: \gene{SK}}{%
(A) Three channel transcripts (\gene{SK}, \gene{Shal}, \gene{slo}) painted on the connectome from each cell's atlas unit:
bright, cell type; dim, hemilineage; dimmer, class; gray, unmapped.
(B) The trace carriers of scene 1 (cells holding the trace at the first decision), by \gene{SK}.
(C) The ventral nerve cord atlas joined to the connectome through \val{M05} hemilineages, ordered by \gene{SK}.
(D) Carrier enrichment in \gene{SK} and in the adaptation index, against matched sets and a mapping-shuffle null, four
learners and the ventral nerve cord alone (open), in null SD; adjusted \Pv{} from the declared families.
(E) Carrier share against \gene{SK} across \val{M05} hemilineages.
(F) Scene-1 carriers mapped by cell type, with hemilineages, and in the ventral nerve cord.
(G) Carrier share and the channel panel across hemilineages: color, z score; size, share of cells detected.
(H) Validation of the bridge: recovered against published cluster sizes, connectome against atlas cells per hemilineage, and
the transmitters and markers that agree.
(I) Carrier types with split-GAL4 lines: each type on the brain, its carriers among its cells, the share of its cue drive into
readout types and its channel strip.
(J) Brains with homogeneous or atlas-informed heterogeneous adaptation, \val{M14} runs each: net memory use joined by run, and
the change when adaptation is removed or states are swapped; bars, medians; exact tests, Holm- and BH-adjusted.
See also Figure~S4.}

\deflegend{5}{Access to the memory is locked to the timing of learning}{%
(A) The cue difference in membrane voltage at the first decision after four delays, with each pattern's angle and cosine to
the trained \val{G06}-ms pattern. Scale bar, 100~μm.
(B) Net memory use on the grid of delay and prefix input, mean of \val{L19} replicates; outlined, trained timings.
(C) The lock (trained minus untrained use) per run in nine scenes, the trained scenes and their variants; bars, medians; exact
sign-flip tests, Holm-adjusted.
(D) Trained against untrained timings, two learners: the adaptation trace at the decision, the cue difference in the readout
input and the push toward the correct target; bars, medians; exact tests, Holm-adjusted.
(E) Each cell's cue difference after \val{G06} against \val{G01}~ms; black, the strongest \val{D15} of cells.
(F) Decision geometry at every timing: the push toward the correct target against the push shared by both cues.
(G) One trace pattern injected at seven delays: the cosine between response directions, and each response's push; shaded,
trained delays.
(H) Training at fixed or jittered timings: untrained use by delay (lines, means of \val{L19} replicates; dots, replicates),
and per replicate inside and outside the jittered range; two-sided sign-flip tests, Holm-adjusted.
See also Figure~S5 and Video~S3.}

\deflegend{6}{The wiring is robust; access follows the trained states}{%
(A) Cells on routes (walks of two or three synapses from the visual input to the readout), intact and after random removal of
\val{G07}, \val{G08} and \val{G02} of the other cells; color, routes through each cell. Scale bar, 100~μm.
(B) Readout cells reachable and routes kept under random or hubs-first removal, and the readout cells by shortest route.
(C) Single extra spikes in five cells at two delays: three kicked cells with their partners, and the cells changed at the
decision.
(D) Mean firing rate against a uniform lowering of the threshold; shaded, runaway.
(E) Training across threshold states (\val{E22} runs) against training at one state (\val{E21}): access window and cell-loss
tolerance relative to the median at one state, and memory under held-out threshold jitter. Bars, geometric means (window,
tolerance) and means (memory); open, beyond the tested range; exact tests, Holm- and BH-adjusted.
(F) The same mechanism in both: the change in net use when adaptation is cleared or states are swapped.
(G) The route network by cell class: strands by transmitter; darker cores, the routes left after \val{G02} removal.
(H) The cumulative share of routes across cell types ranked by traffic, and the cell types carrying routes at each level, by
class.
See also Figure~S6.}

\deflegend{7}{Disease steals access, not memories}{%
(A) The credibility screen of fly pathologies (Table~S1).
(B) Four regimes: the brain state at learning and at test, and the trained histories solved of \val{Q04} (dose 1, readout
trained).
(C) Trained histories solved after a lesion and when learning under it, by dose, two learners; lines, means; Page trend tests,
Holm-adjusted.
(D) Relearning after the lesion: trained histories solved after readout updates, by dose (\val{Q26} runs).
(E) Late restoration of excitability after learning under the lesion (\val{Q25} runs); exact tests, BH-adjusted.
(F) Phenotype profiles of the four credible cases by dose, and where each acts; comparison and equivalence (TOST) tests,
Holm-adjusted.
(G) Disease on the brain: \gene{Shal} transcripts under Aβ42, firing under the uniform Kv4 loss, the Kv4 loss patterned by
\gene{Shal}, and the aging operator. Scale bar, 100~μm.
(H) Task firing rate under the uniform and the \gene{Shal}-patterned Kv4 loss, by dose.
(I) Trained histories solved by learners trained under the uniform lesion: as learned, under the \gene{Shal} pattern and with
excitability restored.
(J) Recall needs the learning state at both ends. Left, trained histories solved when the trace is written in one state and
read in another (the full state moved at the end of the delay), dose 1; means of readouts trained healthy (\val{K18}) or under
the lesion (\val{K19}). Right, across doses, the cosine with the healthy brain of the stored trace and of the healthy trace's
reading (median and range over the trained timings), and the cells holding the trace (median).
See also Figure~S7 and Video~S4.}
